\documentclass[11pt,a4paper]{article}
\usepackage[T1]{fontenc}
\usepackage{lmodern}
\usepackage[margin=25mm,headheight=15pt]{geometry}
\usepackage{amsmath,amssymb,amsthm,mathtools,mathrsfs}
\usepackage{microtype,booktabs,array,longtable,enumitem}
\usepackage{xcolor,fancyhdr,needspace}
\usepackage{tocloft}
\setlength{\cftbeforesecskip}{4pt}
\usepackage{hyperref,bookmark}
\definecolor{ink}{HTML}{183B50}
\hypersetup{colorlinks=true,linkcolor=ink,citecolor=ink,urlcolor=ink,
 pdftitle={Global Fueter Inversion Beyond Finite Connectivity},
 pdfauthor={Research article prepared for Vladimir Reshetnikov},
 pdfsubject={Real-axis gluing, arbitrary polynomial periods, and stable reconstruction}}
\pagestyle{fancy}\fancyhf{}
\fancyhead[L]{\small Global Fueter inversion beyond finite connectivity}
\fancyhead[R]{\small September 2026}
\fancyfoot[C]{\thepage}
\setlength{\parskip}{3pt}\setlength{\parindent}{1.2em}
\setlength{\emergencystretch}{3em}
\numberwithin{equation}{section}
\newtheorem{theorem}{Theorem}[section]
\newtheorem{lemma}[theorem]{Lemma}
\newtheorem{proposition}[theorem]{Proposition}
\newtheorem{corollary}[theorem]{Corollary}
\theoremstyle{definition}
\newtheorem{definition}[theorem]{Definition}
\theoremstyle{remark}
\newtheorem{remark}[theorem]{Remark}
\newcommand{\R}{\mathbb R}
\newcommand{\C}{\mathbb C}
\newcommand{\N}{\mathbb N}
\newcommand{\Z}{\mathbb Z}
\newcommand{\ii}{\mathrm i}
\newcommand{\dd}{\,\mathrm d}
\newcommand{\HH}{\mathbb H}
\newcommand{\Cl}{\operatorname{Cl}}
\newcommand{\Hol}{\mathcal O}
\newcommand{\A}{\mathcal A_h}
\newcommand{\Fu}{\mathcal F_h}
\newcommand{\Pspace}{\mathscr P_h}
\newcommand{\Per}{\operatorname{Per}}
\newcommand{\calW}{\mathscr W}
\newcommand{\Rop}{\mathsf R}
\newcommand{\Sop}{\mathsf S}
\newcommand{\norm}[1]{\left\lVert#1\right\rVert}
\newcommand{\abs}[1]{\left|#1\right|}
\newcommand{\barz}{\overline z}
\newcommand{\dbar}{\partial_{\bar z}}
\newcommand{\wind}{\operatorname{wind}}
\newcommand{\id}{\operatorname{id}}
\newcommand{\rank}{\operatorname{rank}}
\newcommand{\supp}{\operatorname{supp}}
\newcommand{\repoSHA}{cb1646dc442724f8e298cf259195b6c308367d80}

\begin{document}
\hypersetup{pageanchor=false}
\begin{titlepage}
\vspace*{10mm}
\noindent{\large\color{ink} ANALYSIS / TOPOLOGY / RECONSTRUCTION}
\vspace{15mm}
\begin{flushleft}
{\LARGE\bfseries\color{ink} Global Fueter Inversion\\[5pt]
Beyond Finite Connectivity}\par
\vspace{9mm}
{\large Real-axis gluing, arbitrary polynomial periods,\\[3pt]
and the finite--infinite splitting dichotomy}\par
\vspace{13mm}
Research article prepared for Vladimir Reshetnikov\\
30 September 2026
\end{flushleft}
\vspace{7mm}
\begin{abstract}
The polynomial-current formulation of Fueter inversion turns local
holomorphic primitives into a finite collection of closed one-forms. We
extend that formulation to conjugation-symmetric planar domains meeting the
real axis and to domains of arbitrary connectivity. Coalescence of the two
Hermite interpolation nodes is resolved by a contour formula; on the real
axis the interpolant becomes the Taylor polynomial of order $2h-1$. The
resulting current gives a complete global inverse criterion in every even
ambient dimension $2h+2\ge4$.

Every admissible polynomial-valued homology character is realizable. On a
symmetric domain, admissibility means invariance under the homology action of
complex conjugation. Thus a conjugate pair of holes contributes $2h\dim_\R E$
real obstruction coordinates, while a reflection-fixed hole contributes none.
No finiteness or boundary-regularity hypothesis is needed for the general
realization theorem.

The infinite-dimensional problem has a sharp functional-analytic distinction.
The normalized inverse on zero-period targets is continuous linear. By
contrast, realization of arbitrary periods has a continuous linear, or even a
continuous positively homogeneous, right inverse exactly when the admissible
period rank is finite. At infinite rank there is nevertheless a continuous
nonlinear right inverse, and there are continuous linear realization operators
on every prescribed weighted bounded class of period sequences. Complete
proofs, concrete examples, finite symbolic diagnostics, and a research agenda
are provided.
\end{abstract}
\vfill
\noindent\small\textbf{Status and scope.} An unrefereed research manuscript with
written proofs, extending the inspected ProveIt report. Historical priority is
not established, and no Lean or Rocq verification is claimed. Classical
Fueter calculus, Cousin theory, and functional-analytic ingredients are
identified separately from the extensions proved here.
\end{titlepage}
\hypersetup{pageanchor=true}
\pagenumbering{roman}
\setcounter{tocdepth}{2}
{\small\setlength{\parskip}{0pt}\tableofcontents}
\clearpage\pagenumbering{arabic}

\section{The problem, the repository gap, and the results}
\label{sec:intro}

The inverse Fueter problem asks when an axially monogenic function is the
iterated Laplacian of a slice-regular function. Local solvability is not the
same as global single-valued solvability. A local inverse can acquire a
polynomial after continuation around a loop. One must therefore distinguish
three operations: computing the obstruction of a target, constructing a target
with prescribed obstruction, and inverting a target whose obstruction vanishes.
They have different stability properties.

The ProveIt report \emph{Slice regularity and the global inverse Fueter problem}
contains a higher-dimensional construction in its source manuscript 11,
\emph{Polynomial Periods and Sharp Residue Hierarchies for the Global Inverse
Fueter Map} \cite{repo}. That construction applies to connected planar domains
strictly above the real axis. It identifies a polynomial-valued current,
proves its Hermite differential identity, constructs every inverse when its
periods vanish, and realizes all finite collections of period polynomials.
Its further-research questions 1, 4, and 5 ask respectively about crossing the
real axis, stable inverse operators, and infinite connectivity with growth
control. These are the specific starting points of this article.

The inspected repository is pinned to commit
\begin{center}\small\texttt{\repoSHA}.\end{center}
The source report is not formally verified. The present article does not
inherit formal status from unrelated developments in ProveIt. We reprove the
current and its Hermite identity so that the extensions do not rely on an
unexamined higher-dimensional normalization.

\subsection{A summary of the extension}

Fix $h\ge1$ and a nonzero finite-dimensional real orthogonal Clifford module
$E$, of real dimension $d$. Write
\[
 \Pspace=\{M(t):\deg M<2h,\ M\text{ has coefficients in }E\},
 \qquad \dim_\R\Pspace=2hd.
\]
The coefficient field here is real. The symbol $\ii$ in a holomorphic stem is
central and is not a Clifford imaginary unit.

We treat two domain types. A \emph{product domain} is represented by a
connected open $D\subset\{\operatorname{Im}z>0\}$. A \emph{symmetric domain}
is a connected open $D\subset\C$ with $c(D)=D$, where $c(z)=\barz$.
In the latter case stems are intrinsic,
$F(\barz)=\overline{F(z)}$, and the target has the corresponding even/odd
parity. A connected symmetric open set necessarily meets the real axis.

For product domains define
\[
 \calW(D)=\operatorname{Hom}_\R(H_1(D;\R),\Pspace).
\]
For symmetric domains define
\begin{equation}\label{eq:intro-W}
 \calW(D)=\{\rho\in\operatorname{Hom}_\R(H_1(D;\R),\Pspace):
                    \rho\circ c_*=\rho\}.
\end{equation}
Both spaces carry the topology of pointwise convergence on homology classes.

\begin{theorem}[Main conclusions, with notation detailed below]
\label{thm:main}
For either domain type, the following statements hold.
\begin{enumerate}[label=\textup{(\roman*)},leftmargin=2em,itemsep=3pt]
\item The polynomial current in \eqref{eq:current} extends without a pole
through the real axis in the symmetric case. A target has a global Fueter
inverse exactly when all its polynomial periods vanish. The inverse is
obtained by integrating the current and evaluating at its moving complex node.
\item There is an exact sequence of real vector spaces
\begin{equation}\label{eq:intro-exact}
 0\longrightarrow\Pspace\longrightarrow\Hol_*(D;E_\C)
 \xrightarrow{\Fu}\A(D;E)
 \xrightarrow{\Per}\calW(D)\longrightarrow0.
\end{equation}
Every admissible character occurs, including unrestricted infinite period data.
The last map is an open continuous map of Fr\'echet spaces.
\item After a polynomial normalization, the inverse on $\ker\Per$ is
continuous linear, with compact-set estimates valid across the real axis.
\item Let $r$ be the real dimension of $H_1(D;\R)$ in the product case and
of $H_1(D;\R)^+$ in the symmetric case. The period map admits a continuous
linear right inverse if and only if $r<\infty$. The same criterion holds for
a right inverse continuous at zero and positively homogeneous.
\item At $r=\infty$, the period map has a continuous nonlinear right inverse,
but no right inverse can have a continuous linear derivative at even one point.
For each prescribed positive envelope $b=(b_j)$, it has a continuous linear
right inverse on $\ell^\infty(b;\Pspace)$.
\end{enumerate}
\end{theorem}

Here $H_1(D;\R)^+$ is the $+1$ eigenspace of $c_*$. Its dimension is at most
countably infinite. If $r=\infty$, then $\calW(D)$ is a countable
\emph{product} of copies of $\Pspace$, not their direct sum; it does not have
countable algebraic dimension.

The stable-splitting result concerns \emph{realization of arbitrary
obstructions}, not inversion of a compatible target. Confusing those two maps
would reverse the practical meaning of the theorem.

\subsection{What is settled, and what remains outside the claim}

The real-axis question is answered here under smooth axial parity hypotheses,
for every $h$. Infinite connectivity is answered qualitatively without a
boundary assumption, and quantitatively in compact-set Fr\'echet topology
and on prescribed magnitude envelopes. These results do not characterize
Hardy, Bergman, physical $L^p$, or boundary-H\"older period data. They also do
not give sharp constants under degeneration of the domain. Those are genuine
remaining analytic problems, not conclusions hidden in the word
``continuous.''

The forward Fueter theorem and local polynomial kernel are classical
\cite{colombo2011,colombo2014}. The higher-dimensional polynomial current and
its moving Hermite identity are inherited from the repository source and
reproved. The extensions are the confluent real-axis formulation, the complete
arbitrary-domain admissible cokernel, and the resulting realization and
stability classification. The functional-analytic mechanism of nonsplitting
belongs to classical Eidelheit theory \cite{vogt,debrouwere}; we give an
elementary proof and apply it to this precisely identified analytic quotient.
These distinctions matter more than attaching an unsupported priority label
to the whole package.

\section{Spaces, parity, and normalization}
\label{sec:spaces}

\subsection{Coefficients and axial functions}

Put $m=2h+1$. The real Hilbert space $E$ has operators $J_1,\ldots,J_m$
satisfying
\begin{equation}\label{eq:clifford}
 J_j^*=-J_j,\qquad J_jJ_k+J_kJ_j=-2\delta_{jk}\id.
\end{equation}
For $I\in S^{m-1}$ write $J_I=\sum I_jJ_j$ and abbreviate $J_Ia$ by $Ia$.
Then $J_I^2=-\id$ and $J_I$ is an isometry. The complexification
$E_\C=E\oplus\ii E$ is equipped with
\[
 \norm{a+\ii b}^2=\norm a^2+\norm b^2,
 \qquad \overline{a+\ii b}=a-\ii b.
\]
A holomorphic stem $F=A+\ii B$ determines the slice function
$f(x+yI)=A(x,y)+IB(x,y)$. Holomorphicity is
\begin{equation}\label{eq:CR}
 A_x=B_y,\qquad A_y=-B_x.
\end{equation}
In the symmetric case, intrinsic stems have $A(x,-y)=A(x,y)$ and
$B(x,-y)=-B(x,y)$, so this lift is well defined and smooth at the axis.
For $h=1$, $E=\HH$ with left multiplication by the three quaternionic units
is an example. We do not identify the full algebra $\Cl_{0,3}$ with $\HH$.

For a product domain let $\Hol_*(D;E_\C)=\Hol(D;E_\C)$. For a symmetric
domain let
\[
 \Hol_*(D;E_\C)=\{F\in\Hol(D;E_\C):F(\barz)=\overline{F(z)}\}.
\]
These are real vector spaces, with the compact-open topology. Convergence of
all holomorphic derivatives on compact sets follows from Cauchy's estimates.

\subsection{Targets, including the axis}

On a product domain, $\A(D;E)$ consists of smooth pairs $(P,Q)$ solving
\begin{equation}\label{eq:vekua}
 P_x-Q_y=\frac{2h}{y}Q,\qquad P_y+Q_x=0.
\end{equation}
We write $g=P+IQ$ for the physical function and $W=P+\ii Q$ for the pair.
On a symmetric domain we require smooth $P,Q$ on all of $D$, the parity
\begin{equation}\label{eq:parity}
 P(x,-y)=P(x,y),\qquad Q(x,-y)=-Q(x,y),
\end{equation}
and the equations
\begin{equation}\label{eq:smooth-vekua}
 y(P_x-Q_y)=2hQ,\qquad P_y+Q_x=0
\end{equation}
on $D$. These are equations with smooth coefficients, including at $y=0$.
They are equivalent to \eqref{eq:vekua} off the axis, together with smoothness
and parity. The topology on $\A(D;E)$ is the induced $C^\infty$ topology.

The Dirac operator applied away from the axis gives
\[
 \left(\partial_x+\sum_{j=1}^mJ_j\partial_{x_j}\right)(P+IQ)
 =P_x-Q_y-\frac{2h}{y}Q+I(P_y+Q_x).
\]
Thus these are the axial monogenic equations. Our local inverse theorem will
also show directly that every smooth pair admitted at the axis has a genuine
smooth, real-analytic physical lift there. No arbitrary singular solution is
being admitted by multiplying the equation by $y$.

The map under study is the unnormalized differential Fueter map
\begin{equation}\label{eq:Fu}
 \Fu F=\Delta_{2h+2}^{\,h}\mathcal I(F).
\end{equation}
All coefficients and period normalizations below refer to this choice.
Define
\begin{equation}\label{eq:constants}
 c_h=2^hh!,\qquad C_h=(-1)^hc_h^2,\qquad
 \kappa_h=\frac{(-1)^{h+1}}{2^{2h-1}h!(h-1)!}.
\end{equation}
In particular $\kappa_hC_h=-2h$.

\section{The inherited current and its differential identity}
\label{sec:current}

This section rederives the source report's algebraic foundation. It is included
to make the later extension independent of an unverified citation to that
report, not to count the same result as a new contribution.

\subsection{Radial calculus}

On either half-plane put
\[
 \Rop u=y^{-1}\partial_yu,
 \qquad \Sop u=\partial_y(u/y).
\]
These operators apply componentwise.

\begin{lemma}[Radial form and calibration]\label{lem:radial}
For a holomorphic stem $F=A+\ii B$, off the real axis,
\begin{equation}\label{eq:radial}
 \Fu F=P+IQ,\qquad P=c_h\Rop^hA,\qquad Q=c_h\Sop^hB.
\end{equation}
This pair satisfies \eqref{eq:vekua}. For every $a\in E$,
\begin{align}
 \Fu(z^ka)&=0 &&(0\le k<2h),\label{eq:cal-low}\\
 \Fu(z^{2h}a)&=C_ha,\label{eq:cal-even}\\
 \Fu(z^{2h+1}a)&=C_h\bigl((2h+1)xa+Iya\bigr).\label{eq:cal-odd}
\end{align}
\end{lemma}
\begin{proof}
On the scalar axial part the Laplacian is
$L_{2h}=\partial_x^2+\partial_y^2+(2h/y)\partial_y$; on $IB$ it is
$I M_{2h}B$, where $M_s=L_s-s/y^2$. Direct differentiation gives
\[
 L_s\Rop=\Rop L_{s-2},\qquad M_s\Sop=\Sop M_{s-2}.
\]
The planar harmonicity of $A,B$ gives $L_sA=s\Rop A$ and
$M_sB=s\Sop B$. Induction gives a factor
$\prod_{j=0}^{h-1}(2h-2j)=c_h$ after $h$ Laplacians. The identities
\[
 \partial_y\Rop^h=\Sop^h\partial_y,
 \qquad \Rop^h\partial_y=(\partial_y+2h/y)\Sop^h
\]
follow inductively from their definitions. Applying them to \eqref{eq:CR}
proves the target equations.

A slice lift of $z^ka$ is a physical polynomial of degree $k$, so the
order-$2h$ differential operator kills it for $k<2h$. In the real part of
$z^{2h}$ the only surviving highest radial term is $(-1)^hy^{2h}$, and
$\Rop^hy^{2h}=c_h$. The imaginary part is killed. For $z^{2h+1}$ the
surviving terms are $(-1)^h(2h+1)xy^{2h}$ and $(-1)^hy^{2h+1}$; the latter
satisfies $\Sop^hy^{2h+1}=c_hy$. This proves the calibrations.
\end{proof}

\begin{remark}[The real coefficient restriction is essential]
The constant stem $F=\ii a$ has
\[
 \Fu(\ii a)=(-1)^h(2h)!\,Ia/y^{2h},
\]
which is nonzero for $a\ne0$. Even a constant with a purely central-imaginary
coefficient need not lie in the kernel on a product domain. The kernel is
not the space of complex-coefficient polynomials of degree less than $2h$.
\end{remark}

\subsection{A polynomial-valued closed one-form}

Let $t$ be a real formal variable, held fixed under differentiation in $x,y$,
and put $R_t=(t-x)^2+y^2$. Define
\begin{equation}\label{eq:current}
\begin{split}
 \Omega_h[g](t)=\kappa_h R_t^{h-1}\bigl[
 &((x-t)P+yQ)\dd x\\
 &+(yP-(x-t)Q)\dd y\bigr].
\end{split}
\end{equation}
This is a one-form with values in $\Pspace$; it has degree at most $2h-1$ in
$t$. The formula has no negative power of $y$.

\begin{lemma}[Closedness, injectivity, and reflection]\label{lem:current}
For every target in the spaces defined above, $\Omega_h[g]$ is smooth and
closed coefficientwise. The assignment $g\mapsto\Omega_h[g]$ is injective.
On a symmetric domain,
\begin{equation}\label{eq:invariant-current}
 c^*\Omega_h[g]=\Omega_h[g].
\end{equation}
\end{lemma}
\begin{proof}
Off the axis let $a=x-t$, $L=aP+yQ$, and $N=yP-aQ$. The target equations give
\[
 N_x-L_y=2(h-1)Q,\qquad aN-yL=-R_tQ.
\]
The derivative of $R_t^{h-1}$ cancels the first expression when computing
$\partial_x(R_t^{h-1}N)-\partial_y(R_t^{h-1}L)$. For $h=1$ both terms
are zero separately. Thus the current is closed. In the symmetric case its
coefficients are smooth on $D$ by the displayed polynomial formula;
closedness on the dense off-axis set implies closedness everywhere.

At any point with $y\ne0$, evaluation at $t=x$ gives the coefficients
$\kappa_hy^{2h-1}Q$ and $\kappa_hy^{2h-1}P$. Hence a zero current forces
$P=Q=0$ off the axis, and continuity gives the same on the axis.
Finally the $\dd x$ coefficient is even in $y$ and the $\dd y$ coefficient
is odd. Pullback also changes the sign of $\dd y$, proving
\eqref{eq:invariant-current}.
\end{proof}

In central complex notation the same formula is
\begin{equation}\label{eq:complex-current}
\begin{split}
 \Omega_h[g](t)=-\frac{\kappa_h}{2}\bigl[
 &(t-z)^{h-1}(t-\barz)^hW\dd z\\
 &+(t-z)^h(t-\barz)^{h-1}\overline W\dd\barz\bigr].
\end{split}
\end{equation}
Indeed $L-\ii N=(\barz-t)W$ and $L+\ii N=(z-t)\overline W$.
For $h=1$ it is $\theta_g+t\omega_g$, with
\[
 \omega_g=\tfrac12(-P\dd x+Q\dd y),\qquad
 \theta_g=\tfrac12\bigl((xP+yQ)\dd x+(yP-xQ)\dd y\bigr).
\]
Thus the higher-dimensional normalization matches the quaternionic one.

\subsection{Moving Hermite interpolation}

For $z\ne\barz$, let $H_F(z;t)\in\Pspace$ be the unique polynomial satisfying
\begin{align}
 \partial_t^jH_F(z;z)&=F^{(j)}(z),\label{eq:Hermite-jets}\\
 \partial_t^jH_F(z;\barz)&=\overline{F^{(j)}(z)}
                    &&(0\le j<h).\nonumber
\end{align}
This is ordinary two-node Hermite interpolation in each component. Its
coefficients are real: conjugation of all coefficients and interchange of
the two nodes preserves the data, and uniqueness applies.

\begin{theorem}[Hermite differential identity]\label{thm:Hermite}
Off the real axis, for every holomorphic stem,
\begin{equation}\label{eq:Hermite-identity}
 \mathrm d_{x,y}H_F(z;t)=\Omega_h[\Fu F](t).
\end{equation}
\end{theorem}
\begin{proof}
At a fixed $z_0\notin\R$, each side depends real-linearly on the holomorphic
jet through order $h$. Every such jet is the jet of a real-coefficient
polynomial of degree at most $2h+1$: interpolate it and its conjugate through
order $h$ at $z_0,\bar z_0$. The resulting Hermite polynomial has real
coefficients by conjugation and uniqueness. It therefore suffices, for this
arbitrary $h$, to check real monomials of degrees $0,\ldots,2h+1$.

For degrees below $2h$, $H_{z^k}(z;t)=t^k$ and both sides vanish. The remaining
Hermite polynomials are the division remainders
\begin{align}
 H_{z^{2h}}(z;t)&=t^{2h}-R_t^h,\label{eq:rem-even}\\
 H_{z^{2h+1}}(z;t)&=t^{2h+1}-(t+2hx)R_t^h.\label{eq:rem-odd}
\end{align}
They have degree below $2h$ and agree to order $h-1$ at both nodes.
Differentiating the first gives
\[
 -2hR_t^{h-1}\bigl((x-t)\dd x+y\dd y\bigr),
\]
which is the current of \eqref{eq:cal-even} since $\kappa_hC_h=-2h$.
Differentiating the second gives
\[
 -2hR_t^{h-1}\bigl[
 ((2h+1)x(x-t)+y^2)\dd x+y(t+2hx)\dd y\bigr],
\]
which is the current of \eqref{eq:cal-odd}. This proves the identity for every
jet at every point. The argument is an exact finite-dimensional reduction
for arbitrary $h$, not an extrapolation from finitely many tested dimensions.
\end{proof}

\section{Confluence at the real axis and the global inverse}
\label{sec:axis}

The singularity of the usual two-node interpolation matrix at $y=0$ does not
imply a singularity of the interpolant of an intrinsic holomorphic function.
A contour formula makes the cancellation explicit.

\subsection{The confluent interpolation formula}

\begin{theorem}[Regular confluence]\label{thm:confluence}
Let $F\in\Hol_*(D;E_\C)$ on a symmetric domain. The polynomial $H_F$ extends
real-analytically across $D\cap\R$, and its coefficients are even in $y$.
For a real point $x\in D$,
\begin{equation}\label{eq:axis-Taylor}
 H_F(x;t)=\sum_{j=0}^{2h-1}\frac{F^{(j)}(x)}{j!}(t-x)^j.
\end{equation}
The extended polynomial satisfies \eqref{eq:Hermite-identity} on all of $D$.
\end{theorem}
\begin{proof}
Fix a real point $x_0\in D$ and a positively oriented circle $\Gamma$ centered
at $x_0$ with its closed disk in $D$. Restrict $z$ to a smaller disk. Put
\[
 Q_z(\zeta)=((\zeta-z)(\zeta-\barz))^h.
\]
Define a polynomial in $t$ by
\begin{equation}\label{eq:contour-H}
 H(z;t)=\frac1{2\pi\ii}\int_\Gamma
 F(\zeta)\,
 \frac{Q_z(\zeta)-Q_z(t)}{(\zeta-t)Q_z(\zeta)}\dd\zeta.
\end{equation}
The numerator divided by $\zeta-t$ is polynomial in $t$ of degree $2h-1$.
The other denominator has no zero on $\Gamma$. Therefore each coefficient
of $H$ is real-analytic in $(x,y)$, including $y=0$, and is even in $y$.

For $t$ inside $\Gamma$, Cauchy's formula rewrites this as
\[
 H(z;t)=F(t)-Q_z(t)\frac1{2\pi\ii}
          \int_\Gamma\frac{F(\zeta)}{(\zeta-t)Q_z(\zeta)}\dd\zeta.
\]
At $y\ne0$ the last term vanishes to order $h$ at each node. Intrinsicness
identifies the data at $\barz$ with the conjugate data at $z$, so uniqueness
gives $H=H_F$. In particular its coefficients are in $E$, also at the axis
by continuity.

At $y=0$, $Q_x(t)=(t-x)^{2h}$, and the same expression shows that $H(x;t)$
interpolates the jet of $F$ through order $2h-1$ at $x$. Its degree bound
proves \eqref{eq:axis-Taylor}. The local extensions agree off the axis and
hence on overlaps. Finally both sides of \eqref{eq:Hermite-identity} are
smooth at the axis, and they agree on its complement, so they agree everywhere.
\end{proof}

\begin{remark}
The current itself needed no new regularizing factor: it was already
polynomial in $y$. What needed regularization was the representation of its
primitive by a two-node interpolation matrix. Formula \eqref{eq:contour-H}
resolves exactly that degeneration. It is a confluent Hermite formula, not
a new claim about Cauchy integration.
\end{remark}

\subsection{Exactness, reconstruction, and every inverse}

For a piecewise smooth closed curve $\gamma$ define
\begin{equation}\label{eq:period}
 \Per_g(\gamma;t)=\int_\gamma\Omega_h[g](t)\in\Pspace.
\end{equation}
Closedness makes this depend only on homology. We extend real-linearly to
$H_1(D;\R)$. In the symmetric case \eqref{eq:invariant-current} implies
$\Per_g\circ c_*=\Per_g$.

\begin{theorem}[Complete inverse criterion on both domain types]
\label{thm:inverse}
A target $g\in\A(D;E)$ has a global inverse in $\Hol_*(D;E_\C)$ if and only
if $\Per_g=0$. When this holds, choose a basepoint $z_0\in D$, real in the
symmetric case, and any $H_0\in\Pspace$. Then
\begin{equation}\label{eq:inverse-integral}
 H(z;t)=H_0(t)+\int_{z_0}^{z}\Omega_h[g](t),
 \qquad F(z)=H(z;z)
\end{equation}
is an inverse. The integral is path independent. Every inverse arises in this
way, and any two inverses differ by a real polynomial in $\Pspace$.
\end{theorem}
\begin{proof}
If $g=\Fu F$, the Hermite identity, including its confluent version, makes
$\Omega_h[g]$ exact. Thus all periods vanish.

Conversely zero periods imply coefficientwise path independence in
\eqref{eq:inverse-integral}; connected open planar sets are path connected.
The resulting smooth real-coefficient polynomial satisfies
$\mathrm dH=\Omega_h[g]$. In the symmetric case, invariance of the current
and a real basepoint imply $H(\barz;t)=H(z;t)$, coefficientwise: conjugate a
path from $z_0$ to $z$ and use \eqref{eq:invariant-current}. It follows that
$F(\barz)=\overline{F(z)}$.

Formula \eqref{eq:complex-current} yields
\begin{align}
 \partial_zH(z;t)&=-\frac{\kappa_h}{2}
             (t-z)^{h-1}(t-\barz)^hW,\label{eq:Hz}\\
 \dbar H(z;t)&=-\frac{\kappa_h}{2}
             (t-z)^h(t-\barz)^{h-1}\overline W.\label{eq:Hbarz}
\end{align}
In the derivative $\dbar H(z;z)$ the substituted variable $t=z$ contributes
zero, and the right side of \eqref{eq:Hbarz} vanishes at $t=z$. Thus $F$ is
holomorphic, including at the axis.

Off the axis we claim
$F^{(j)}(z)=\partial_t^jH(z;z)$ for $0\le j<h$. The case $j=0$ is the
definition. If $j\le h-2$, differentiating in $z$ adds
$\partial_z\partial_t^jH(z;t)|_{t=z}$, which is zero by the factor
$(t-z)^{h-1}$ in \eqref{eq:Hz}. This proves the induction. Real coefficients
supply the conjugate-node data. Hence $H=H_F$ off the axis by interpolation
uniqueness. Consequently
\[
 \Omega_h[\Fu F]=\mathrm dH_F=\mathrm dH=\Omega_h[g].
\]
Injectivity of the current gives $\Fu F=g$, with equality on the axis by
continuity. The holomorphic intrinsic stem provides its smooth physical
lift there.

For exhaustiveness, every inverse has a Hermite polynomial whose differential
is the same current. Their coefficientwise difference is constant on connected
$D$, so it belongs to $\Pspace$. Evaluation at $t=z$ gives precisely the
stated polynomial ambiguity. Conversely the calibrations kill every such
real polynomial.
\end{proof}

\begin{corollary}[Local solvability, kernel, and normalization]
\label{cor:local}
Every target has a local inverse, also near a real point. Its local
coefficients $P,Q$ are real-analytic. On either connected domain type,
\[
 \ker\Fu=\Pspace.
\]
Taking $H_0=0$ in \eqref{eq:inverse-integral} fixes a unique global inverse
whenever periods vanish. On a product domain this normalization is
$F^{(j)}(z_0)=0$ for $0\le j<h$. On a symmetric domain it is
$F^{(j)}(x_0)=0$ for $0\le j<2h$.
\end{corollary}
\begin{proof}
Use a small off-axis disk, or a small symmetric disk centered at the real
point. Every closed form is exact there; in the latter case use a real
basepoint. The constructed inverse is holomorphic. Its radial image is
real-analytic off the axis. At a real center its convergent intrinsic power
series lifts to a convergent physical power series, whose iterated Laplacian
is real-analytic. This proves the analytic assertion on all of $D$.

The kernel follows from Theorem~\ref{thm:inverse} applied to zero. The
normalizations are equivalent to $H_F(z_0;t)=0$ by the two-node jet conditions
or the confluent Taylor formula, respectively.
\end{proof}

The change from $h$ complex jet conditions to $2h$ real jet conditions at a
real basepoint preserves exactly $2hd$ real constraints. It is not a loss of
half of the polynomial gauge.

\section{Realizing every admissible period character}
\label{sec:realization}

Finite logarithmic sums are not enough to justify arbitrary infinite period
assignments: an uncontrolled infinite sum can diverge on every compact set.
We first prove unrestricted existence by a local-to-global argument, and only
then address continuous choices and controlled sums.

\subsection{The classical inputs}

We use three standard facts in precise forms. First, de Rham's theorem in
degree one identifies smooth closed real one-forms modulo exact forms on an
open planar domain with $\operatorname{Hom}_\R(H_1(D;\R),\R)$, via integration.
It applies componentwise to any finite-dimensional real coefficient space.
Second, the additive Cousin problem is solvable on every open planar domain:
a holomorphic additive one-cocycle is a difference of local holomorphic
functions. Equivalently, every smooth equation $\dbar v=\beta$ with smooth
coefficient $\beta$ is globally solvable there. These are ordinary open
Riemann-surface results; no compact support or boundary regularity is imposed
\cite{forster}. Third, the open mapping theorem applies to continuous linear
surjections between Fr\'echet spaces. We will use that theorem only after
proving surjectivity and identifying both topologies.

For clarity, the Cousin step used below can be reduced explicitly to the
stated $\bar\partial$ theorem. Suppose $c_{ij}$ is a holomorphic additive
cocycle on a locally finite cover and $(\chi_k)$ is a subordinate smooth
partition of unity. Put
\begin{equation}\label{eq:Cousin-partition}
 u_i=\sum_k\chi_k c_{ik}.
\end{equation}
Each term is extended by zero where necessary, using the support of $\chi_k$.
The cocycle equation gives $u_i-u_j=c_{ij}$. Hence $\dbar u_i$ is a global
smooth coefficient $\beta$. Choose $v$ with $\dbar v=\beta$ and set
$F_i=u_i-v$. Then $F_i$ is holomorphic and $F_i-F_j=c_{ij}$.
This is the only global holomorphic existence input in our period realization.

\subsection{Unrestricted realization theorem}

\begin{theorem}[Arbitrary admissible periods]\label{thm:surjectivity}
For either domain type, every $\rho\in\calW(D)$ is $\Per_g$ for some
$g\in\A(D;E)$. In particular, \eqref{eq:intro-exact} is exact without a
connectivity bound or a boundary assumption.
\end{theorem}
\begin{proof}
By de Rham's theorem choose a smooth closed $\Pspace$-valued one-form $\alpha$
whose periods are $\rho$. In the symmetric case replace it by
$(\alpha+c^*\alpha)/2$. Its periods remain $\rho$, since
$\rho\circ c_*=\rho$, and now $c^*\alpha=\alpha$.

Take a locally finite cover by sufficiently small open disks. In the symmetric
case take a cover permuted by conjugation, using paired disks disjoint from the
axis and fixed disks centered on the real axis. Such a cover and a subordinate
partition of unity can be chosen equivariantly: first take a locally finite
refinement of small disks and its reflection, then average the associated
partition under the two-element group. Intersections of the chosen disks are
convex whenever nonempty.

On each disk $U_i$ choose a real-polynomial-valued primitive $H_i$ of
$\alpha$. In the symmetric case index the conjugate disk by $ci$ and choose
\begin{equation}\label{eq:primitive-reflection}
 H_{ci}(\barz;t)=H_i(z;t).
\end{equation}
For a pair this is a definition of the second primitive. On a fixed disk it
follows by choosing a real basepoint: the two sides have the same differential
and the same value there.

On an overlap, the difference
\[
 C_{ij}(t)=H_i(z;t)-H_j(z;t)
\]
is a locally constant element of $\Pspace$. Evaluation gives a holomorphic
additive cocycle $c_{ij}(z)=C_{ij}(z)$. Solve the additive Cousin problem to
obtain holomorphic $F_i$ with
\begin{equation}\label{eq:cocycle-F}
 F_i-F_j=C_{ij}(z).
\end{equation}

In the symmetric case the solution can be chosen equivariantly. Indeed choose
the partition in \eqref{eq:Cousin-partition} so that
$\chi_{ci}(\barz)=\chi_i(z)$. Real coefficients and
\eqref{eq:primitive-reflection} then give
$u_{ci}(\barz)=\overline{u_i(z)}$. For a smooth function let
$v^\sharp(z)=\overline{v(\barz)}$. Direct differentiation shows
\[
 \dbar v^\sharp(z)=\overline{(\dbar v)(\barz)}.
\]
The global coefficient $\beta$ has this same symmetry. Therefore replacing
any solution $v$ by $(v+v^\sharp)/2$ retains $\dbar v=\beta$ and gives
$F_{ci}(\barz)=\overline{F_i(z)}$. In particular $F_i$ is intrinsic on every
fixed disk.

The local targets $\Fu F_i$ agree on overlaps by \eqref{eq:cocycle-F} and the
real polynomial kernel. They define a smooth target $g$ on all of $D$; in
the symmetric case it has the required parity and is smooth at every real
point. On each disk its Hermite polynomial satisfies
\[
 \mathrm dH_{F_i}=\Omega_h[g].
\]
The difference $H_{F_i}-H_i$ is globally defined: on an overlap,
\[
 H_{F_i}-H_{F_j}=H_{C_{ij}(z)}=C_{ij}(t)=H_i-H_j.
\]
This identity holds also at the axis by confluence. Hence for a global smooth
polynomial-valued function $B$,
\[
 \Omega_h[g]-\alpha=\mathrm dB.
\]
Their periods agree, so $\Per_g=\rho$. The equality also fixes the sign of
the prescribed monodromy without a convention-dependent Cech calculation.
The other exactness assertions were proved in Theorem~\ref{thm:inverse}.
\end{proof}

\begin{remark}[Why no additional real-axis obstruction appears]
The reflection group is finite and the coefficient space is real of
characteristic zero, so the averaging operations above are legitimate.
Intrinsicness near a real point is secured on a fixed disk before applying
the Fueter map. One is not gluing arbitrary off-axis stems across the axis
and hoping that their values match. Invariance of the character is both
necessary and sufficient.
\end{remark}

\subsection{The actual size of the obstruction space}

Open planar domains have the homotopy type of countable CW complexes, so
$H_1(D;\R)$ has a finite or countable algebraic basis; this is also a
consequence of countable topology for open Riemann surfaces \cite{forster}.
In the symmetric case $c_*^2=1$, and the operators
\[
 \pi_+=\tfrac12(1+c_*),\qquad \pi_-=\tfrac12(1-c_*)
\]
split the real homology into its two eigenspaces. An invariant character kills
$\operatorname{im}\pi_-$ and is arbitrary on $\operatorname{im}\pi_+$.
Thus, after choosing a basis of the relevant homology space,
\begin{equation}\label{eq:W-product}
 \calW(D)\cong
 \begin{cases}
  \Pspace^{\,r},&r<\infty,\\
  \Pspace^{\N},&r=\infty.
 \end{cases}
\end{equation}
The topology is the product topology: each homology vector is a finite linear
combination of basis vectors, so pointwise evaluation and coordinatewise
convergence give the same topology.

In particular finite rank gives real cokernel dimension $2hdr$. At infinite
rank, infinitely supported arbitrary coordinate sequences occur. There is no
hidden summability condition in unrestricted smooth axial analysis. Growth
conditions enter only after a stronger analytic space has been specified.

\section{Reflection-fixed holes, conjugate pairs, and examples}
\label{sec:holes}

\subsection{Finite connectivity with a geometric basis}

Assume here that $D$ is bounded, connected, symmetric, and its boundary is a
finite union of disjoint $C^1$ Jordan curves. Suppose its bounded complementary
components consist of $p$ conjugate pairs and $q$ reflection-fixed components.
Thus the total number of holes is $b=2p+q$.

\begin{theorem}[Which holes obstruct inversion]\label{thm:holes}
Choose counterclockwise generators $\gamma_j^+,\gamma_j^-$ about the two
members of each conjugate pair, and $\delta_k$ about each reflection-fixed
hole. Then every smooth target satisfies
\begin{equation}\label{eq:paired-periods}
 \Per_g(\gamma_j^-)=-\Per_g(\gamma_j^+),
 \qquad \Per_g(\delta_k)=0.
\end{equation}
The polynomials $\Per_g(\gamma_j^+)$ are otherwise arbitrary. Therefore
\begin{equation}\label{eq:finite-dimension}
 \dim_\R\operatorname{coker}\Fu=2hdp.
\end{equation}
\end{theorem}
\begin{proof}
Conjugation reverses planar orientation, and consequently
\[
 c_*[\gamma_j^+]=-[\gamma_j^-],\qquad
 c_*[\gamma_j^-]=-[\gamma_j^+],\qquad
 c_*[\delta_k]=-[\delta_k].
\]
The invariance condition on the character proves \eqref{eq:paired-periods}.
Each paired block contributes one $+1$ eigenvector; a fixed block contributes
none. For example $([\gamma_j^+]-[\gamma_j^-])/2$ is a basis of the invariant
part with character value $\Per_g(\gamma_j^+)$. Its dimension is $p$.
Theorem~\ref{thm:surjectivity} gives arbitrary values on these vectors,
proving the dimension formula.
\end{proof}

Notice that the character is \emph{invariant under the homology involution},
while its values on the two \emph{counterclockwise} generators are opposite.
These are the same condition, not a choice between two sign conventions.

\subsection{An explicit paired logarithmic representative}

For a point $a\notin D$, with $\bar a\notin D$, and $M\in\Pspace$, consider
local branches of
\begin{equation}\label{eq:paired-log}
 F_{a,M}(z)=\frac{M(z)}{2\pi\ii}
               \operatorname{Log}\frac{z-a}{z-\bar a}.
\end{equation}
Equivalently one may use the difference of two local logarithms. Different
branches differ by an integer multiple of the real polynomial $M(z)$, so
$G_{a,M}=\Fu F_{a,M}$ is single valued. Near a real point the ratio has modulus
one; a local logarithm can be chosen purely imaginary on the real interval.
Then $F_{a,M}$ is intrinsic there. More generally conjugation changes a
branch to the conjugate intrinsic branch modulo an integer multiple of $M$.
Thus the target is a smooth symmetric axial target everywhere on $D$.

Analytic continuation and the Hermite identity give
\begin{equation}\label{eq:paired-wind}
 \Per_{G_{a,M}}(\gamma;t)=
       \bigl(\wind(\gamma,a)-\wind(\gamma,\bar a)\bigr)M(t).
\end{equation}
For finitely many conjugate pairs, choose a point in one member of each pair
and sum these targets. This is an explicit linear right inverse of the finite
period map. For a product domain one instead uses
$M(z)\operatorname{Log}(z-a)/(2\pi\ii)$, as in the source report.

The targets need not be rational functions, nor need all logarithmic terms
vanish from a radial expression. What vanishes is the Fueter image of the
\emph{branch difference}. That is sufficient, and is the exact reason that
the construction patches.

\subsection{Examples distinguishing topology from admissible rank}

\paragraph{A real annulus.}
Let $D=\{r<|z|<R\}$ with $0<r<R$. Its first Betti number is one, but its
single hole is reflection fixed. Hence every smooth intrinsic axial target
has a global Fueter inverse, for every $h\ge1$. The same statement holds on
$\C\setminus\{0\}$. The inverse may be singular at the omitted real point;
surjectivity is not a removability theorem.

\paragraph{A pair of punctures.}
Let $D=\C\setminus\{\ii,-\ii\}$. The admissible period rank is one, so the
real cokernel dimension is $2hd$. The target associated with
\[
 F(z)=\frac{M(z)}{2\pi\ii}\operatorname{Log}\frac{z-\ii}{z+\ii}
\]
has upper period $M$ and lower period $-M$. It has a global single-valued
inverse precisely when $M=0$. For $h=1$, $E=\HH$, this is an
$8$-dimensional real obstruction; for general $h$ the factor is $2hd$.

\paragraph{Infinitely many real punctures.}
Let $D=\C\setminus\Z$. Compact cycles meet only finitely many of the
punctures in a bounded region, and the homology is generated by the small
counterclockwise puncture loops. Conjugation acts as minus the identity on
each generator. Thus $r=0$ although the first homology has countably infinite
rank. Every target has a global normalized inverse, continuously and linearly.
This example rules out a criterion stated merely as ``finitely many holes
versus infinitely many holes.'' The correct invariant is admissible period
rank.

\paragraph{Infinitely many conjugate pairs.}
Let
\[
 D=\C\setminus\{2^j+\ii,2^j-\ii:j\ge1\}.
\]
Now $r=\infty$. Any sequence $(M_j)\in\Pspace^\N$, however rapidly its
coefficients grow, is realized by a smooth axial target. A naive series of
the paired logarithmic targets need not converge. Sections~\ref{sec:nonsplit}
and~\ref{sec:selection} show both why no universal continuous linear remedy
exists on the product space and what positive replacements are possible.
The product domain $\{\operatorname{Im}z>0\}\setminus\{2^j+\ii:j\ge1\}$
has the corresponding unrestricted countable obstruction space without a
reflection condition.

\section{The topological exact sequence and a stable compatible inverse}
\label{sec:topology}

\subsection{Fr\'echet spaces and a continuous norm}

\begin{lemma}[Topologies and unique continuation]\label{lem:Frechet}
The spaces $\Hol_*(D;E_\C)$ and $\A(D;E)$ are real Fr\'echet spaces, and
$\Fu$ is continuous linear. For any closed disk $K\Subset D\setminus\R$
with nonempty interior, the function
\begin{equation}\label{eq:continuous-norm}
 q_K(g)=\sup_{z\in K}\norm{P(z)+\ii Q(z)}
\end{equation}
is a continuous norm on $\A(D;E)$. It is not asserted to generate the
Fr\'echet topology.
\end{lemma}
\begin{proof}
The holomorphic space with compact-open topology is Fr\'echet; its intrinsic
part is closed under the continuous real involution $F\mapsto F^\sharp$.
The target space is a closed subspace of $C^\infty(D;E^2)$: the differential
equations have smooth coefficients in the product case, and
\eqref{eq:smooth-vekua} and the parity equations are closed conditions in the
symmetric case. Hence it is Fr\'echet.

Continuity of $\Fu$ off the axis follows from the radial formula and Cauchy
estimates. Near a real center $x_0$, expand an intrinsic stem as
$F(z)=\sum_{n\ge0}a_n(z-x_0)^n$ with $a_n\in E$. On any smaller disk,
Cauchy estimates bound its coefficients by a fixed geometric sequence times
a supremum norm on a larger disk. Its physical lift is the corresponding
convergent series of homogeneous slice polynomials. Any fixed finite number
of derivatives multiplies the degree-$n$ estimates by at most a polynomial
in $n$ and reduces the power of the radius by a fixed amount. The resulting
geometrically convergent majorant bounds the iterated Laplacian and every
fixed target derivative. A finite cover of a compact set proves continuity
also across the axis.

The expression $q_K$ is a continuous seminorm. If it is zero, both $P,Q$
vanish on an open disk. By Corollary~\ref{cor:local} these coefficients are
real-analytic throughout the connected set $D$. The identity theorem for
real-analytic functions gives $P=Q=0$ on $D$. Thus $q_K$ is a norm.
\end{proof}

\begin{proposition}[The period space is the quotient topology]
\label{prop:quotient}
The period map is continuous and open onto $\calW(D)$. Its kernel is closed,
and
\begin{equation}\label{eq:quotient}
 \A(D;E)/\operatorname{im}\Fu\ \cong\ \calW(D)
\end{equation}
is an isomorphism of topological real vector spaces. Also
$\Hol_*(D;E_\C)/\Pspace\cong\ker\Per$ topologically.
\end{proposition}
\begin{proof}
Each coordinate of a period character is a finite linear combination of
integrals over fixed compact curves. The current has smooth polynomial
coefficients, so this coordinate is continuous in the target topology.
The product topology \eqref{eq:W-product} is exactly the topology determined
by these coordinates. Theorem~\ref{thm:surjectivity} gives surjectivity, and
the open mapping theorem for Fr\'echet spaces makes the map open. Its kernel
is closed since the product target is Hausdorff. Theorem~\ref{thm:inverse}
identifies it with $\operatorname{im}\Fu$, proving \eqref{eq:quotient}.
The same open mapping argument applies to the induced continuous bijection
from the Fr\'echet quotient $\Hol_*/\Pspace$ to the closed Fr\'echet subspace
$\ker\Per$.
\end{proof}

This proposition prevents an artificial topology from creating the later
nonsplitting phenomenon. The product topology is the natural quotient
topology of the smooth target space itself.

\subsection{A useful direct period estimate}

For real $t$ and a parametrized curve with unit tangent $(u,v)$, the integrand
inside the current is
\[
 P((x-t)u+yv)+Q(yu-(x-t)v).
\]
The two scalar coefficients have squared sum
$((x-t)^2+y^2)(u^2+v^2)$. Cauchy--Schwarz therefore gives
\begin{equation}\label{eq:period-bound}
 \norm{\Per_g(\gamma;t)}
 \le |\kappa_h|\int_\gamma
       ((t-x)^2+y^2)^{h-1/2}\norm W\dd s.
\end{equation}
This supplies a direct continuity estimate and a finite-curve error estimate
for each real evaluation of the polynomial. Values at $2h$ distinct real
$t$ determine its coefficients by a fixed finite-dimensional interpolation
map. Small numerical periods, however, do not certify exactly zero periods.

\subsection{Compact-set estimates for the normalized inverse}

\begin{theorem}[Stable inversion on compatible targets]\label{thm:stable-inverse}
Fix the normalization of Corollary~\ref{cor:local}. The resulting map
\[
 \mathcal R:\ker\Per\longrightarrow\Hol_*(D;E_\C)
\]
is continuous linear and satisfies $\Fu\mathcal Rg=g$. More explicitly, for
any compact $K\Subset D$ and integer $s\ge0$, there are a compact $L\Subset D$
and $C>0$, depending on $K,s,D,h$ and the fixed basepoint, such that
\begin{equation}\label{eq:inverse-compact-estimate}
 \sup_{z\in K}\norm{(\mathcal Rg)^{(s)}(z)}
       \le C\sup_{w\in L}\norm{W(w)}
       \qquad(g\in\ker\Per).
\end{equation}
In a symmetric domain $K$ and the integration paths may cross the real axis;
no factor $1/\min|\operatorname{Im}w|$ is needed.
\end{theorem}
\begin{proof}
Linearity follows from the path integral with $H_0=0$. First consider a compact
set $K'$ with a neighborhood in $D$. Cover it by finitely many small disks
with compact closures in $D$ and join their centers to the basepoint by
finitely many piecewise smooth paths. Any point of $K'$ can be reached by
one of these paths followed by a straight segment. Thus all paths lie in a
fixed compact $L$, and their lengths are bounded by a constant $\ell$.

Choose $R\ge1$ bounding the moduli of all endpoints and all path points.
In the integral formula for $F(z)$, evaluate the formal variable at the
fixed endpoint $t=z$. At a path point $w=x+\ii y$,
\[
 |x-z|\le2R,\quad |y|\le R,\quad
 |(z-x)^2+y^2|\le5R^2.
\]
The triangle inequality and Cauchy--Schwarz bound the bracket in
\eqref{eq:current}, evaluated at $t=z$, by $6R\norm W\dd s$.
Consequently
\begin{equation}\label{eq:explicit-path-bound}
 \sup_{z\in K'}\norm{F(z)}
 \le 6|\kappa_h|R(5R^2)^{h-1}\ell\sup_{w\in L}\norm{W(w)}.
\end{equation}
This bound uses no division by a path's imaginary coordinate.

For general $s$, take $K'$ large enough to contain closed disks of some common
radius $\delta>0$ around every point of $K$, all inside $D$. The Cauchy
estimate multiplies \eqref{eq:explicit-path-bound} by $s!\delta^{-s}$.
This proves \eqref{eq:inverse-compact-estimate} and continuity in the
compact-open holomorphic topology.
\end{proof}

The theorem supplies a local, axis-regular answer to the source report's
stability question. It does not imply a bounded operator between any specified
global Sobolev spaces: the compact $L$ and path geometry are allowed to depend
on $K$. Nor does it provide a continuous linear way to remove arbitrary
periods from an incompatible target. The latter operation is addressed next.

\section{The finite--infinite obstruction to linear splitting}
\label{sec:nonsplit}

A right inverse of the period map is a choice of a target for each admissible
period character. This is a different right inverse from the normalized
Fueter inverse $\mathcal R$ of the previous section.

\subsection{An elementary product-space obstruction}

\begin{lemma}[Continuous norms exclude homogeneous sections]
\label{lem:no-section}
Let $V\ne0$ be a finite-dimensional real vector space, let $X$ be a real
topological vector space with a continuous norm $q$, and let
$T:X\to V^\N$ be linear. There is no map $S:V^\N\to X$ satisfying all three
conditions
\[
 TS=\id,\qquad S(\lambda v)=\lambda S(v)\ (\lambda\ge0),
 \qquad S\text{ is continuous at }0.
\]
In particular, no continuous linear right inverse exists. Every continuous
linear map $L:V^\N\to X$ factors through finitely many coordinates.
\end{lemma}
\begin{proof}
Homogeneity implies $S(0)=0$. Continuity at zero gives a basic product
neighborhood $U$ such that $q(S(v))<1$ for $v\in U$. Some finite set of
coordinates determines $U$; increasing its largest index if necessary, take
that set inside $\{1,\ldots,N\}$. Choose a nonzero sequence $v$ vanishing in
these coordinates. Then $\lambda v\in U$ for every $\lambda>0$, so
\[
 \lambda q(S(v))=q(S(\lambda v))<1\qquad(\lambda>0).
\]
Therefore $q(S(v))=0$, which means $S(v)=0$, since $q$ is a norm. This
contradicts $TS(v)=v\ne0$.

For a continuous linear $L$ the identical argument shows that it vanishes
on every sequence whose first $N$ coordinates are zero. Thus it depends only
on those coordinates, and has finite-dimensional range.
\end{proof}

The linear mechanism is classical in the theory of surjections onto the
sequence space $\omega=\R^\N$; compare the continuous-norm obstruction in
Vogt's Proposition 3.1 \cite{vogt} and the vector-valued Eidelheit setting
\cite{debrouwere}. Its application here is not a claim to a new abstract
nonsplitting theorem. The additional point is that the analytic period
quotient has now been proved to be exactly such a product space, including
its correct real-axis symmetry.

\subsection{The sharp criterion for Fueter period realization}

\begin{theorem}[Splitting dichotomy]\label{thm:split}
Let $r$ be the admissible period rank in Theorem~\ref{thm:main}. The following
are equivalent:
\begin{enumerate}[label=\textup{(\alph*)},leftmargin=2em,itemsep=2pt]
\item $r<\infty$.
\item $\Per:\A(D;E)\to\calW(D)$ has a continuous linear right inverse.
\item It has a positively homogeneous right inverse continuous at zero.
\item $\ker\Per$ is a complemented closed subspace of $\A(D;E)$.
\end{enumerate}
If $r=\infty$, a section cannot have, at any point, a directional derivative
in every direction which is a continuous linear operator. In particular it
cannot be a $C^1$ section in the usual locally convex sense.
\end{theorem}
\begin{proof}
If $r<\infty$, choose a basis of the finite-dimensional period space and a
target realizing each basis vector, using Theorem~\ref{thm:surjectivity}.
Linear extension gives a continuous section because its domain is finite
dimensional. For regular finite-connectivity domains the paired logarithmic
formula gives such a choice explicitly. This proves (a)$\Rightarrow$(b),
and (b)$\Rightarrow$(c) is immediate.

If $r=\infty$, identify $\calW(D)$ with $\Pspace^\N$ and apply
Lemma~\ref{lem:no-section}, using the continuous norm from
Lemma~\ref{lem:Frechet}. Hence (c)$\Rightarrow$(a).
A continuous linear section $S$ gives the continuous projection
$\id-S\Per$ onto $\ker\Per$. Conversely, if a continuous projection $P$
onto that kernel exists, its complementary range $X_0=(\id-P)\A(D;E)$ is a
closed Fr\'echet subspace. The restriction $\Per:X_0\to\calW(D)$ is a
continuous linear bijection, and its inverse is continuous by the open
mapping theorem. This proves equivalence with (d).

Finally suppose a section has the stated derivative $L$ at $v_0$. For every
direction $v$, apply the continuous linear map $\Per$ to the difference
quotient:
\[
 \Per\!\left(\frac{S(v_0+\varepsilon v)-S(v_0)}{\varepsilon}\right)=v.
\]
Taking the limit yields $\Per L=\id$. This is a continuous linear right
inverse, which is impossible at infinite rank.
\end{proof}

\begin{corollary}[No universal linear removal of infinitely many periods]
\label{cor:no-removal}
At infinite admissible period rank there is no continuous linear projection
which replaces an arbitrary target by a zero-period target while fixing all
zero-period targets. Nevertheless the normalized inverse of each zero-period
target remains continuous linear.
\end{corollary}
\begin{proof}
The hypothetical projection would complement $\ker\Per$, contradicting the
previous theorem. The positive statement is Theorem~\ref{thm:stable-inverse}.
\end{proof}

The product topology controls any chosen finite collection of periods but
places no bound on the remaining ones. The impossibility proof exploits
exactly those uncontrolled amplitudes. It is not an argument based on a
particular logarithmic series, on a poorly chosen basis, or on numerical
ill-conditioning of a matrix.

\section{Nonlinear and envelope-controlled sections}
\label{sec:selection}

Nonsplitting does not mean that continuous realization is impossible. The
following constructive-in-form argument identifies a continuous nonlinear
choice and linear choices on stronger input spaces. The existence of the
local lifts used in the construction is supplied by the open mapping theorem;
we do not claim an effective algorithm for finding those lifts on an arbitrary
domain.

\subsection{A small-coordinate lifting lemma}

\begin{lemma}[Arbitrarily small lifts of sufficiently late coordinates]
\label{lem:small-lifts}
Let $T:X\to V^\N$ be a continuous linear surjection from a Fr\'echet space,
where $V$ is finite dimensional. Choose increasing seminorms
$q_1\le q_2\le\cdots$ defining the topology and put $q_0=0$. There are
strictly increasing positive integers $m(n)\ge n$ such that, for every
$j>m(n)$ and $\varepsilon>0$, there exists a linear map $A:V\to X$ with
\begin{equation}\label{eq:small-lift}
 TA(v)=e_jv,\qquad q_n(A(v))\le\varepsilon\norm v.
\end{equation}
Here $e_jv$ denotes the sequence supported at coordinate $j$ with value $v$.
For $n=0$ one may take $m(0)=0$ and any exact linear coordinate lift.
\end{lemma}
\begin{proof}
By openness, $T(\{x:q_n(x)<1\})$ contains a basic product neighborhood whose
restricted coordinates are among the first $m(n)$. Enlarge these indices so
that they are strictly increasing and $m(n)\ge n$. If $j>m(n)$, every scalar
multiple $\lambda e_jv$ lies in that neighborhood. It therefore has a lift
$x_\lambda$ with $q_n(x_\lambda)<1$. Dividing by $\lambda>0$ gives an exact
lift of $e_jv$ with arbitrarily small $q_n$.

Take a fixed orthonormal basis $v_1,\ldots,v_s$ of $V$ and lift each $e_jv_a$
with $q_n$ at most $\varepsilon/\sqrt{s}$. Linear extension gives
\[
 q_n\!\left(\sum_a\lambda_a A(v_a)\right)
 \le \frac{\varepsilon}{\sqrt{s}}\sum_a|\lambda_a|
 \le\varepsilon\norm{\sum_a\lambda_av_a}.
\]
Each such linear map is continuous because its domain is finite dimensional.
For $q_0=0$ only existence of the finitely many lifts is needed.
\end{proof}

Define
\begin{equation}\label{eq:nj}
 n(j)=\max\bigl(\{0\}\cup\{n\ge1:m(n)<j\}\bigr).
\end{equation}
This is finite for every $j$, and $n(j)\to\infty$. The finitely many early
coordinates with $n(j)=0$ create no convergence problem.

\subsection{A continuous nonlinear realization of all sequences}

\begin{theorem}[Continuous nonlinear section]\label{thm:nonlinear}
Every continuous linear surjection $T:X\to V^\N$ as above has a continuous
right inverse $S$ with $S(0)=0$. It can be written as
\begin{equation}\label{eq:nonlinear-series}
 S((v_j))=\sum_{j=1}^\infty\Phi_j(v_j),
\end{equation}
where each $\Phi_j:V\to X$ is continuous, $T\Phi_j(v)=e_jv$, and for every
$n$ and $N\ge m(n)$,
\begin{equation}\label{eq:nonlinear-tail}
 q_n\!\left(S((v_j))-\sum_{j=1}^N\Phi_j(v_j)\right)\le2^{-N}
\end{equation}
for all input sequences, without a bound on their amplitudes.
\end{theorem}
\begin{proof}
For every $j\ge1$ and integer $k\ge0$, use
Lemma~\ref{lem:small-lifts} to choose a linear coordinate lift $A_{j,k}:V\to X$
satisfying
\begin{equation}\label{eq:Ajk}
 TA_{j,k}(v)=e_jv,
 \qquad q_{n(j)}(A_{j,k}(v))\le\frac{2^{-j}}{k+1}\norm v.
\end{equation}
For a vector $v$, put $r=\norm v$, $k=\lfloor r\rfloor$, and
$\theta=r-k\in[0,1)$. Define
\begin{equation}\label{eq:Phi}
 \Phi_j(v)=(1-\theta)A_{j,k}(v)+\theta A_{j,k+1}(v).
\end{equation}
At an integer radius the two adjacent definitions meet at the same value
$A_{j,k}(v)$, so this map is continuous. It vanishes at zero and has exact
coordinate image. Since $r\le k+1$ and $k+2\ge k+1$, the triangle inequality
in \eqref{eq:Ajk} gives
\begin{equation}\label{eq:uniform-small}
 q_{n(j)}(\Phi_j(v))\le2^{-j}\qquad(v\in V).
\end{equation}

For a fixed seminorm $q_n$, all $j>m(n)$ satisfy $n(j)\ge n$. The tail of
\eqref{eq:nonlinear-series} is therefore absolutely summable in $q_n$,
uniformly over every input sequence. The finitely many earlier terms are
well defined. Completeness of $X$ gives convergence in its Fr\'echet topology,
and summing the geometric bound gives \eqref{eq:nonlinear-tail}.

Each partial sum is continuous in the product topology because it depends
on finitely many coordinates. Uniform convergence in each defining seminorm
makes the limit continuous: to control a finite collection of seminorms,
first choose a common tail cutoff, then use continuity of the finite partial
sum. Applying the continuous $T$ to the limit gives each coordinate exactly,
so $TS=\id$. Finally $S(0)=0$ term by term.
\end{proof}

For $T=\Per$, the theorem applies after
Proposition~\ref{prop:quotient}. At infinite rank the amplitude-dependent
choice of $A_{j,k}$ is indispensable: keeping a single lift per coordinate
would give a linear section forbidden by Theorem~\ref{thm:split}.
The uniform smallness of late coordinates is only a compact-seminorm
statement; their realizations can become arbitrarily large elsewhere in the
domain.

\begin{corollary}[A nonlinear topological splitting]
\label{cor:homeomorphism}
The map
\[
 g\longmapsto\bigl(g-S(\Per_g),\Per_g\bigr)
\]
is a homeomorphism
\begin{equation}\label{eq:homeomorphism}
 \A(D;E)\ \cong\ \ker\Per\times\calW(D).
\end{equation}
Its inverse sends $(k,\rho)$ to $k+S(\rho)$. At infinite admissible rank this
splitting cannot be linear or positively homogeneous with a continuous
realization section.
\end{corollary}
\begin{proof}
The first component has zero periods since $\Per S=\id$. The two displayed
maps are continuous and mutually inverse. The impossibility statement is
Theorem~\ref{thm:split}.
\end{proof}

Combining the homeomorphism with $\mathcal R$ gives a continuous nonlinear
decomposition into a normalized stem and an obstruction representative. It
is a topological splitting, not a Fr\'echet linear splitting in disguise.

\subsection{Linear realization on every prescribed magnitude envelope}

Fix any positive sequence $b=(b_j)_{j\ge1}$ and define the Banach space
\begin{equation}\label{eq:weighted-space}
 \ell^\infty(b;V)=
 \left\{v=(v_j):\norm v_b:=\sup_j\frac{\norm{v_j}}{b_j}<\infty\right\}.
\end{equation}
No growth bound is imposed on the chosen envelope itself. It can increase or
decrease arbitrarily rapidly.

\begin{theorem}[Envelope-controlled linear sections]\label{thm:weighted}
For every positive envelope $b$, a continuous linear surjection
$T:X\to V^\N$ as above has a continuous linear right inverse
\[
 S_b:\ell^\infty(b;V)\longrightarrow X.
\]
It can be chosen as an exact-coordinate series
$S_b(v)=\sum_{j\ge1}B_j(v_j)$ such that, for $N\ge m(n)$,
\begin{equation}\label{eq:weighted-tail}
 q_n\!\left(S_b(v)-\sum_{j=1}^N B_j(v_j)\right)
       \le2^{-N}\norm v_b.
\end{equation}
On each bounded coefficient box $\{v:\norm v_b\le R\}$ the section is also
continuous for the product topology, and its image is compact in $X$.
\end{theorem}
\begin{proof}
Choose linear coordinate lifts $B_j$ with
\[
 TB_j(v)=e_jv,\qquad
 q_{n(j)}(B_j(v))\le\frac{2^{-j}}{b_j}\norm v,
\]
using Lemma~\ref{lem:small-lifts}. For any fixed $n$, sufficiently late terms
satisfy $q_n(B_j(v_j))\le2^{-j}\norm v_b$. Completeness gives convergence,
linearity is preserved by the series, and geometric summation yields
\eqref{eq:weighted-tail}. Each finite initial sum is bounded from
$\ell^\infty(b;V)$ into each seminorm of $X$, since every $B_j$ is a
finite-dimensional continuous map and each $b_j$ is finite. Adding the
uniform tail bound proves continuity. Exactness holds coordinatewise under $T$.

On a fixed bounded coefficient box the same tail bound is uniform over all
inputs. Finite partial sums are product-continuous, so the limit is
product-continuous on that box. The box is the product of finite-dimensional
closed bounded balls and is compact. Its continuous image in $X$ is compact.
\end{proof}

Applied to Fueter periods, the theorem supplies a convergent linear splitting
for any \emph{fixed} coefficient magnitude class, with seminorm tail estimates.
Every individual sequence belongs to some such class, for example with
$b_j=1+\norm{v_j}$. But there is no single continuous linear operator on the
whole product which simultaneously implements these choices. The quantifier
``for each envelope there exists a section'' cannot be interchanged with
``there exists a section for all envelopes.''

These are not sharp physical-growth estimates near an accumulating boundary.
The integers $m(n)$ and the coordinate lifts are obtained existentially from
openness. Effective bounds in terms of conformal geometry, and membership in
specified global analytic norm spaces, remain separate questions.

\section{Worked formulas, diagnostics, and reproducibility}
\label{sec:verification}

\subsection{A six-dimensional confluent example}

Take $h=2$, so $c_2=8$, $C_2=64$, and $\kappa_2=-1/16$. For the intrinsic
stem $F(z)=z^4a$, the target is the constant $64a$ and its Hermite polynomial is
\begin{equation}\label{eq:example-H}
 H_F(z;t)=\bigl[t^4-((t-x)^2+y^2)^2\bigr]a.
\end{equation}
At $y=0$ this is
\[
 \bigl[x^4+4x^3(t-x)+6x^2(t-x)^2+4x(t-x)^3\bigr]a,
\]
exactly the cubic Taylor polynomial of $z^4a$ at $x$. The nonzero target is
encoded by the variation of that polynomial with its basepoint, not by a
singularity of its coefficients at the axis.

At the real basepoint $x_0=0$, $H_F(0;t)=0$, so $F(z)=z^4a$ is already the
normalized inverse of $64a$. More generally the normalized inverse of a
constant target $a$ on a domain containing a real basepoint $x_0$ is
$(z-x_0)^{2h}a/C_h$.

\subsection{A period check with three loops}

For $h=1,2,3$, use the scalar coefficient direction of a fixed nonzero vector
in $E$ and the local stem
\begin{equation}\label{eq:numerical-stem}
 F(z)=\frac{1+z}{2\pi\ii}
             \bigl(\operatorname{Log}(z-\ii)
                         -\operatorname{Log}(z+\ii)\bigr).
\end{equation}
The degree-one polynomial $1+t$ belongs to $\Pspace$ for every such $h$.
The predicted period polynomials on counterclockwise circles about
$\ii,-\ii,0$ of radii $1/5,1/5,1/4$, respectively, are
\begin{equation}\label{eq:test-periods}
 1+t,\qquad -(1+t),\qquad 0.
\end{equation}
The third loop crosses the real axis. The first two test the orientation sign;
the third tests cancellation in a region with no admissible hole.

The code computes holomorphic jets of \eqref{eq:numerical-stem} explicitly,
forms the radial derivatives, and integrates every coefficient of the target
current by periodic trapezoidal quadrature. Differentiating a discontinuous
numerically chosen logarithm branch is avoided: only analytic local jet
formulas are used. Real-polynomial branch changes are annihilated by the
forward map. The quadrature nodes are shifted away from the exact real-axis
points when using the off-axis radial expression; the theorem itself has no
such exclusion.

\subsection{What the executed tests check}

The supplied \texttt{code/verify.py} was executed using SymPy 1.14.0 and
mpmath 1.3.0. The retained JSON and text log are in \texttt{data/}.
The exact tests cover closedness and normalization for $1\le h\le8$,
reflection invariance of both current coefficients, the two injectivity
evaluations, and the nonzero Fueter image of the central-imaginary constant.
They also verify the Hermite differential identity and real-axis confluence
on real monomials for $1\le h\le5$ and degrees $0\le k\le2h+4$, including
all calibration monomials through degree $2h+1$.

The homology test checks the reflection matrix with paired negative-swap
blocks and fixed negative-identity blocks for every $0\le p,q\le6$. Its
$+1$ eigenspace has dimension $p$ in all 49 cases. A symbolic identity also
checks the scalar coefficient norm used in \eqref{eq:period-bound}.

The numerical part uses 60 decimal digits and 192 nodes for each of nine
contour tests: three values of $h$ and the three loops just described.
The largest observed absolute coefficient error over these nine tests was
below $2.7\times10^{-49}$. The complete errors are in the log. These are
finite-precision observations, not certified quadrature enclosures. The global topology,
Cousin realization, infinite-rank impossibility, and infinite-series convergence
are proved in the text, not tested by a finite computation.

\subsection{Dependency and status audit}

\begin{center}
\small
\begin{tabular}{>{\raggedright\arraybackslash}p{0.29\textwidth}
 >{\raggedright\arraybackslash}p{0.32\textwidth}
 >{\raggedright\arraybackslash}p{0.27\textwidth}}
\toprule
Statement & Proof mechanism & Status in this article\\
\midrule
Radial Fueter formula and kernel & Radial commutators and calibration & Classical, reproved\\
Polynomial current and Hermite identity & Target equations; universal jet reduction & Repository foundation, reproved\\
Axis confluence and inverse & Contour Hermite formula; exact integration & Extension proved\\
Arbitrary admissible cokernel & de Rham; additive Cousin; reflection averaging & Extension proved using stated classical inputs\\
Compact inverse estimates & Bounded path networks; Cauchy estimates & Extension proved\\
Finite--infinite splitting criterion & Continuous norm; product topology & Application of classical mechanism, proved\\
Nonlinear and envelope sections & Small coordinate lifts; uniformly controlled series & Explicit abstract constructions, proved\\
Boundary-norm classification & Not supplied by compact seminorm estimates & Remains open here\\
\bottomrule
\end{tabular}
\end{center}

No theorem in the article is marked as formally checked. A passed symbolic
test is not a Lean proof, and an all-order proof is not validated merely by
agreement in a finite set of dimensions. Conversely, the infinite statements
do not rest on unperformed infinite computations: their proofs are the
Cousin construction, the product-neighborhood argument, and the seminorm
convergence estimates.

\section{Relation to the literature and to the source questions}
\label{sec:literature}

The inverse Fueter literature contains explicit local and integral inverse
constructions. Colombo, Sabadini, and Sommen develop an inverse mapping
theorem \cite{colombo2011}; Colombo, Pe\~na Pe\~na, Sabadini, and Sommen give
radial integral formulas \cite{colombo2014}. These supply the classical
forward/inverse context. The present paper does not claim to discover local
surjectivity or a polynomial kernel.

For comparison with global quaternionic statements, Perotti's Theorem 2
assumes that every connected component of the planar stem domain is simply
connected \cite{perotti}. Our nontrivial-period examples do not contradict
that result. We retain its domain hypothesis when making the comparison,
rather than treating a broad introductory description of surjectivity as an
unconditional theorem.

There is also a substantial global slice-regular gluing context. Prezelj and
Vlacci construct quaternionic Cartan coverings suited to symmetry and Cousin
problems \cite{prezelj}. Our proof uses only additive, finite-dimensional
coefficient cocycles and a real averaging argument; it does not solve a new
noncommutative multiplicative Cousin problem. The smooth global
$\bar\partial$ theorem and the degree-one de Rham correspondence are standard
inputs from open Riemann-surface theory \cite{forster}.

On the functional-analytic side, surjections onto countable sequence products
are Eidelheit problems. Vogt explicitly uses the fact that a Fr\'echet space
with a continuous norm cannot have a continuous linear section of a quotient
onto $\omega$ \cite{vogt}. Debrouwere and Vindas study vector-valued Eidelheit
sequences in a broader setting \cite{debrouwere}. The benefit here is to
identify the \emph{correct Fueter period quotient} first, including symmetry
and topology, and then separate exact inversion, linear realization, nonlinear
realization, and weighted realization.

A recent related preprint of De Martino and Pinton studies variations of
inverse Fueter constructions and generalized polyanalytic
Cauchy--Kovalevskaya extension \cite{demartino2026}. It is part of the active
inverse-map context, not a premise of the proofs in this article. A focused
inspection of these sources and the pinned repository does not establish
historical priority for the theorem package presented here. The defensible
claim is the collection of explicitly proved statements, not a declaration
that no analogous statement could be found elsewhere.

Relative to the source report's own research list, the outcome is precise.
Question 1, real-axis crossing, is answered by
Theorems~\ref{thm:confluence}, \ref{thm:inverse}, \ref{thm:surjectivity}, and
\ref{thm:holes}. Question 4 receives compact-set inverse estimates and a
sharp obstruction to universal linear period removal; boundary Sobolev and
H\"older estimates remain open here. Question 5 receives unrestricted
realization, continuous nonlinear realization, and linear realization on
any chosen coefficient envelope; sharp analytic growth conditions near
accumulating boundary components remain open here. No broader resolution is
implied by these scoped answers.

\section{Further research questions and topics}
\label{sec:future}

The following questions are not assumptions in any proof above. They separate
several next steps that should not be conflated into a single request for
``better estimates.''

\subsection{Period images in fixed analytic norm spaces}

Fix a multiply connected domain and a specified target norm, for example a
weighted planar $L^2$ norm of $P+\ii Q$ compatible with physical volume.
Characterize the image of its finite-norm monogenic subspace under $\Per$.
For infinitely many holes, is that image describable by a capacity-weighted
sequence space or a non-diagonal quadratic form? Which geometric separation
conditions make a diagonal description equivalent to the true norm?

The envelope theorem is not such a characterization: it constructs smooth
representatives controlled on compact sets, not representatives with a
prescribed global integrability bound. On punctured domains the source
report's residue-energy obstruction already excludes many nonzero periods
from finite-energy classes. Any proposed characterization must respect those
local obstructions before making global summability claims.

\subsection{Effective coordinate lifts and certified truncation}

The proofs produce thresholds $m(n)$ and coordinate lifts with arbitrarily
small specified seminorms, but do not compute them from the geometry of $D$.
For domains with an effective exhaustion and computable boundary curves,
construct a Runge--Cousin algorithm that outputs these lifts with certified
bounds. Then \eqref{eq:weighted-tail} would become a numerical stopping
criterion rather than an existential estimate.

A useful first model is the punctured half-plane with punctures
$2^j+\ii$. One can try to subtract finite holomorphic Taylor or Runge
corrections from each logarithmic stem so that its Fueter image is small on
successive compact sets. The correction degree must depend on the required
amplitude bound; a fixed universal linear correction scheme on all product
data is ruled out by Theorem~\ref{thm:split}.

\subsection{Collapsing conjugate holes and real-axis limits}

A conjugate pair contributes $2hd$ obstruction coordinates, while a
reflection-fixed limiting hole contributes none. Analyze a family of domains
in which a conjugate pair approaches the real axis and merges. What rate of
blow-up is forced on targets whose nonzero upper period is kept fixed? Which
rescalings of the period polynomial have nontrivial limits?

The regularity of the current at $y=0$ on a fixed domain does not make such a
domain-degeneration limit uniform. The geometry of the paths and the local
singularities change simultaneously. Combining the present symmetry
classification with the source report's residue filtration appears a natural
way to formulate the correct transition problem.

\subsection{Boundary estimates and optimal derivative gain}

For finitely connected domains with controlled smooth boundaries, determine
sharp estimates for the normalized inverse in specified Sobolev or H\"older
scales. Identify the role of the polynomial normalization, boundary traces,
and the distance between boundary components. Near the real axis one should
use the confluent formula rather than accept artificial factors arising from
an ill-conditioned two-node interpolation matrix.

A proof should specify the exact domain and range norms and the derivative
gain. The order of the ambient Laplacian alone does not automatically settle
the appropriate gain on the constrained holomorphic/axial subspaces.

\subsection{Least-energy period representatives}

For a finite set of periods and a Hilbert norm on admissible targets, ask
whether a unique norm-minimizing representative exists, modulo the relevant
closed zero-period subspace. Determine its period matrix and geometric
asymptotics. For infinitely many periods, identify when the resulting
minimization problem remains well posed and what sequence space its data must
belong to.

This would give a canonical choice within a normed analytic class. It cannot
produce a continuous linear representative for completely unrestricted
product data at infinite admissible rank, because that would contradict the
splitting theorem.

\subsection{The best regularity of nonlinear sections}

At infinite rank there are continuous sections but no section with a
continuous linear derivative at any point. Determine whether natural sections
can have useful quantitative moduli of continuity in explicit compatible
metrics, or be locally Lipschitz on specified weighted subsets. How do those
moduli depend on the exhaustion and on the magnitude envelope?

The amplitude interpolation in \eqref{eq:Phi} is simple and guarantees
continuity, but it does not optimize regularity or symmetry. A proposed
improvement must not accidentally restore positive homogeneity on the whole
product space, which is impossible.

\subsection{Angular modes and nonaxial reconstruction}

The article fixes the axial mode and a coefficient Clifford module. Extend
the current-and-gluing method to higher spherical monogenic modes, with the
correct radial normalization and angular multiplicities. Determine whether
modewise admissible periods can be assembled into a single obstruction
space for a convergent nonaxial expansion.

Independent formal inverses in each angular mode are not sufficient: one
needs convergence, compatibility at the axis, and a common analytic norm in
which the reconstruction is defined. These issues are distinct from the
coefficientwise Cousin argument used here.

\subsection{A staged formalization and symbolic interface}

A first formalization target is the universal finite-jet reduction proving
Theorem~\ref{thm:Hermite}, followed by the confluent remainder formula and
closedness of the current. These algebraic/differential components are much
smaller than a complete open-Riemann-surface library. The next interface is
exact coefficientwise integration versus vanishing periods; the final one is
an abstract surjective Fr\'echet map into a countable product.

For a symbolic or numerical implementation, record polynomial coefficient
conventions, loop orientations, reflection relations, and the normalization
of the Laplacian explicitly. Numerical estimates of small periods should be
reported as intervals or error bounds, not silently replaced by exact zero.
A useful verified output would be either an exact zero-period certificate or
a certified obstruction, followed by an inverse only in the former case.

\section{Conclusion}

The real axis is not a singularity of the polynomial obstruction current.
The apparent singularity lies in a nonconfluent interpolation representation,
and the contour formula replaces it by a Taylor polynomial. This yields a
complete global inverse criterion with the correct parity and a transparent
reflection action on periods.

The obstruction space is exactly the space of admissible polynomial homology
characters, on arbitrary planar domains. In symmetric geometry a fixed hole
contributes no obstruction and a conjugate pair contributes one polynomial.
The resulting rank, rather than the raw number of holes, governs the
functional-analytic behavior of period realization.

Compatible inversion is continuously linear after a finite polynomial gauge
choice. Realization of unrestricted periods is continuously linearly split
only at finite admissible rank. At infinite rank, continuous nonlinear and
fixed-envelope linear splittings still exist, with controlled seminorm tails.
This is a genuine distinction between topology, existence, and stable choice,
and it identifies which parts of the repository's proposed global extension
can be completed without imposing additional boundary norms.

\appendix
\section{Proof audit: hypotheses that cannot be dropped silently}
\label{app:audit}

\paragraph{Connectedness.}
A single disk norm detects the entire target only because the domain is
connected and the coefficients satisfy unique continuation. On a disconnected
domain the kernel, normalization, and product-space argument must be treated
componentwise; a seminorm on one component is not a norm on all targets.

\paragraph{Nonzero coefficient module.}
The nonsplitting theorem assumes $E\ne0$, so the coordinate space $\Pspace$
is nonzero. For the zero module every space is zero, and the infinite-rank
obstruction statement would be vacuous rather than true as phrased.

\paragraph{Smooth parity at the axis.}
The symmetric target space imposes smoothness on all of $D$ and the specified
even/odd parity. A pair solving the off-axis equations but blowing up as
$y\to0$ is not a target in this theorem. An arbitrary holomorphic stem on the
upper half-plane is not automatically intrinsic across the axis.

\paragraph{Real rather than complex polynomial coefficients.}
The current and its periods take values in $\Pspace$ over $\R$. This is needed
for branch differences to lie in the kernel, for reflection averaging, and
for the precise real dimension. The central-imaginary constant test in the
code is a direct guard against the wrong coefficient field.

\paragraph{Orientation.}
The invariant-character condition is $\rho\circ c_*=\rho$. A
counterclockwise loop reflects to a clockwise loop; hence upper and lower
counterclockwise periods are negatives. The factor $1/2$ in a convenient
invariant homology basis prevents a second, independent factor-of-two error.

\paragraph{Ordinary, not compactly supported, cohomology.}
The realization theorem concerns ordinary periods and unrestricted smooth
targets. The de Rham and Cousin steps do not impose compact support or a
boundary norm. Adding such a condition changes the problem and may change
which characters are realizable.

\paragraph{Product rather than direct-sum data.}
Every unrestricted sequence is allowed. This is crucial both to arbitrary
period realization and to the no-homogeneous-section proof. The latter uses
arbitrarily large tails inside a basic product neighborhood, not a bounded
set in a Banach sequence norm.

\paragraph{What the positive tail estimates mean.}
The indices $m(n)$ are chosen from openness, not supplied as explicit geometric
constants. The formulas establish convergence and continuity, and they give
an exact geometric tail bound once those choices are made. They do not by
themselves constitute an effective numerical solver on a general domain.

\paragraph{Which facts are imported.}
The main global proof imports degree-one de Rham theory, global additive
Cousin/$\bar\partial$ solvability on open planar domains, and the Fr\'echet
open mapping theorem. No global inverse Fueter theorem is imported into the
proof of our inverse criterion. No infinite-dimensional selection theorem is
needed for the explicit section construction of Section~\ref{sec:selection}.

\section{Reproducing the artifact}

The package contains the present \texttt{article.tex} and \texttt{article.pdf},
a shell build script, the exact/numerical diagnostic program, its recorded
output, a concise proof-audit file, and source-provenance notes. The bibliography
is internal to the TeX source. There are no external figures or font files.
From the package directory run
\begin{verbatim}
python code/verify.py
./build.sh
\end{verbatim}
The diagnostic program accepts \texttt{--skip-numerical} and
\texttt{--max-h 1} through \texttt{--max-h 8}. Its default is the retained run:
monomial tests through $h=5$ and current tests through $h=8$. The build script
uses three pdfLaTeX passes to resolve cross-references and the contents.

\begin{thebibliography}{99}
\bibitem{repo}
V. Reshetnikov, \emph{ProveIt}, research collection,
\emph{Slice regularity and the global inverse Fueter problem}; in particular
source 11, \emph{Polynomial Periods and Sharp Residue Hierarchies for the Global
Inverse Fueter Map}, September 2026.
Inspected at commit \texttt{cb1646dc442724f8e298cf259195b6c308367d80}.
\href{https://github.com/VladimirReshetnikov/ProveIt/blob/cb1646dc442724f8e298cf259195b6c308367d80/SetTheory/Cardinals/docs/reports/quaternionic-analysis/slice-regularity-and-fueter-inversion/sources/11-polynomial-periods.tex}{Commit-pinned source manuscript}.
Unrefereed repository research report.

\bibitem{colombo2011}
F. Colombo, I. Sabadini, and F. Sommen,
\emph{The inverse Fueter mapping theorem},
Communications on Pure and Applied Analysis \textbf{10} (2011), 1165--1181.
\href{https://doi.org/10.3934/cpaa.2011.10.1165}{doi:10.3934/cpaa.2011.10.1165}.

\bibitem{colombo2014}
F. Colombo, D. Pe\~na Pe\~na, I. Sabadini, and F. Sommen,
\emph{A new integral formula for the inverse Fueter mapping theorem},
Journal of Mathematical Analysis and Applications \textbf{417} (2014), 112--122.
\href{https://arxiv.org/abs/1302.0685}{arXiv:1302.0685}.

\bibitem{perotti}
A. Perotti, \emph{A local Cauchy integral formula for slice-regular functions},
Computational Methods and Function Theory \textbf{24} (2024), 185--203.
\href{https://doi.org/10.1007/s40315-023-00485-5}{doi:10.1007/s40315-023-00485-5}.
In particular, Theorem 2 and its simply connected component hypothesis.

\bibitem{forster}
O. Forster, \emph{Lectures on Riemann Surfaces}, Graduate Texts in Mathematics,
vol.~81, Springer, 1981; corrected printing 1999.
\href{https://doi.org/10.1007/978-1-4612-5961-9}{doi:10.1007/978-1-4612-5961-9}.
See the treatment of noncompact Riemann surfaces, countable topology,
Mittag--Leffler theory, and prescribed summands of automorphy.

\bibitem{prezelj}
J. Prezelj and F. Vlacci, \emph{Quaternionic Cartan coverings and applications},
The Journal of Geometric Analysis \textbf{35} (2025), article 80.
\href{https://doi.org/10.1007/s12220-025-01900-0}{doi:10.1007/s12220-025-01900-0}.

\bibitem{vogt}
D. Vogt, \emph{On two problems of Mityagin},
Mathematische Nachrichten \textbf{141} (1989), 13--25.
\href{https://www2.math.uni-wuppertal.de/~vogt/preprints/Eigene/44.pdf}{Author-hosted article}.
See especially Proposition 3.1.

\bibitem{debrouwere}
A. Debrouwere and J. Vindas,
\emph{Vector-valued Eidelheit sequences and the non B-completeness of tensor
products}, preprint, 2016.
\href{https://arxiv.org/abs/1612.06560}{arXiv:1612.06560}.

\bibitem{demartino2026}
A. De Martino and S. Pinton,
\emph{A variation of the inverse Fueter theorem and the generalized polyanalytic
Cauchy--Kovalevskaya extension of order 2}, preprint, 7 August 2026.
\href{https://arxiv.org/abs/2608.07711}{arXiv:2608.07711}.

\end{thebibliography}
\end{document}
